%! Rational Functions and Vertical Asymptotes — Unit 3, Lesson 3.3 — Warm-Up (Answer Key)
\documentclass[10pt]{article}
\usepackage{precalculus-article}
\usepackage{precalculus-key}
\newcommand{\TallMath}[1]{$\displaystyle #1\rule[-1.4em]{0pt}{3.2em}$}

\begin{document}

\pageheader{Unit 3, Lesson 3.3}{Warm-Up --- Answer Key}
\namedateperiod

\vspace{0.14in}

\begin{enumerate}[itemsep=14pt]

\item Let $r(x) = \dfrac{x^2 - 4}{x^2 - x - 6}$.

\begin{enumerate}[label=(\alph*), itemsep=6pt]
  \item Factor the numerator completely. \quad \ans{$(x-2)(x+2)$}

  \item Factor the denominator completely. \quad \ans{$(x-3)(x+2)$}

  \item List the zeros of the \textbf{numerator}: \ans{$x = 2,\; x = -2$} \quad
        List the zeros of the \textbf{denominator}: \ans{$x = 3,\; x = -2$}

  \item Do any zeros of the denominator also appear in the numerator? Circle one:
        \quad \ans{Yes} \quad / \quad \textbf{No}

        \textit{(The value $x = -2$ is a zero of both the numerator and denominator.)}
\end{enumerate}

\item For $r(x) = \dfrac{x + 1}{x^2 - 9}$, for what value(s) of $x$ is $r(x)$ \textbf{undefined}?

\vspace{0.1in}
\ans{$x^2 - 9 = (x-3)(x+3) = 0$, so $r(x)$ is undefined at $x = 3$ and $x = -3$.}

\item The table below shows values of $r(x)$ as $x$ approaches 2 from the right.

\smallskip
\renewcommand{\arraystretch}{1.3}
\begin{tabular}{|c|c|c|c|}
\hline
$x$    & 2.1 & 2.01 & 2.001 \\
\hline
$r(x)$ & 10  & 100  & 1{,}000 \\
\hline
\end{tabular}
\smallskip

Based on the table, describe what appears to be happening to $r(x)$ as $x$ approaches 2
from the right:

\vspace{0.1in}
\ans{$r(x)$ is growing larger and larger without bound --- it is increasing without limit.}

\vspace{0.1in}
What do you predict $r(x)$ approaches? \ans{$+\infty$ (positive infinity)}

\end{enumerate}

\end{document}
